\section{Discussion}\label{sec:discussion}

\subsection{Why the network does not help}\label{sec:why-not}

The measurements give three reasons, in decreasing order of how directly they are measured. First, the iteration the network feeds forgets its start. At the published freeze index the run of the iteration passes through eleven applications of the update before the focal phase is frozen (Proposition~\ref{prop:washout}), and the phases of a converged run, the best start that was constructed, save \OTwelveGain{} iteration on ten geometries and one seed (Table~\ref{tab:oracle}), which is consistent with a small value of a phase predictor at $\nfix=12$ and is not a bound over all starts, and the network at $\nfix=12$ saves none (Table~\ref{tab:ablation}).
Second, where the start does matter, at an early freeze, what matters is whether the frozen phases are compatible with a uniform solution, and whether the trap list determines such phases is not tested here: the oracle start, whose converged phases the trap list does not supply, converges on geometries where a random start does not and halves the median at $\nfix=3$ (\ORThreeMed{} to \OThreeMed{}), while the network, trained on a loss that rewards low spread after ten iterations, does not (Table~\ref{tab:sweep}, Table~\ref{tab:oracle}). That headroom of a factor of two is the open question that this network leaves. Third, the training objective did not transfer to the held-out geometries: during training the network lowered the spread after the first iteration
from \trainSigRandFirst{} to \trainSigNetFirst{} percent on the training distribution, an advantage that is gone by iteration ten (\trainSigNet{} against \trainSigRand{} percent) and absent after the first iteration on the held-out geometries (\trNThreeOne{} against \trRTwelveOne{} percent). We cannot say from these runs whether a longer or differently designed training would move the third reason; the first two do not depend on the training.

\subsection{The gain belongs to the freeze index}\label{sec:freeze-gain}

The reduction in median iterations that the candidate shows with respect to the default is real and is reproduced without a network: a random start with the same early freeze needs \itRThree{} iterations. The sweep of Section~\ref{sec:abl-sweep} puts it in context: the median iterations fall steadily as the freeze index falls from 12 to 4 and the share of runs that never converge rises at the same time, from \swTwelveNever{} to \swFourNever{} percent. An early freeze is a bet that the phases at that moment are good
enough, and on this model the bet fails on a share of the geometries that grows as the freeze comes earlier. A practitioner who wants fewer iterations on arrays of the type of the held-out set can therefore lower $\nfix$ and accept the failure share; we did not test any remedy for the failed runs, such as a restart with a later freeze, and we make no recommendation. The default of Kim et al.~\cite{kim2019} was chosen on their own array and camera; we did not test their array, and nothing in these results says that their choice was suboptimal for it. What the results say
is that on the dense generated arrays of this study, in the scalar model, the default is a conservative point of a trade between iterations and reliability, and a learned start does not move the trade.

\subsection{Cost of the network against what it can save}\label{sec:cost}

By operation count (Section~\ref{sec:complexity}) the network pass costs of the order of two iterations of the point form at the largest array of 1024 traps and less on smaller ones, because its attention is quadratic in $M$ while an iteration is linear in it. A start that saved about one iteration at $\nfix=12$ would not repay a pass of that size on the largest arrays, and for arrays of a hundred traps the pass is much cheaper than an iteration, so the net balance is a question of the array size and of the saving, which is zero here. We did not measure the inference time of the exported graph in a browser, so the cost statement is an operation count, not a timing.

The balance is different for the many holograms of a sequence, where each hologram starts from its predecessor and no start has to be predicted, and for arrays smaller than a hundred traps, where the iteration is cheap and the median iterations to 1 percent are already few. The cost statement therefore holds for the dense, large arrays for which a warm start was proposed, which are also the arrays where the pass is the most expensive.

\subsection{The explanation by the weights}\label{sec:error}

The hypothesis that a warm start should be hard because the weights restart at one is not supported by the accounting of Section~\ref{sec:abl-weights}: the allowance permits the weights to rebuild their spread in \wTwelveKMed{} iterations, far before 1 percent arrives. The iteration count is set by the freeze, not by the rebuilding of the weights, in the data of this study. The hypothesis follows from the structure of the iteration and was tested by the measurement of Table~\ref{tab:weights}, which did not support it. The proposition is a true inequality and stays; the hypothesis is dropped. An explanation that follows from the structure of an algorithm is a hypothesis, and the quantity it predicts (here a lower bound on the iterations) can be computed on every run at no cost; computing it shows that the bound is far below the observed counts. The same accounting applied to the oracle start gives the same answer, so the conclusion does not depend on the quality of the start. What the accounting leaves open is why the freeze index sets the count, for which the evidence is that of Conjecture~\ref{conj:freeze}, with its counterexamples, and no proof.

\subsection{What the result does not close}\label{sec:open}

A label-based warm start (regressing converged phases on a family of geometries where they are determined, for instance regular lattices) or a modified freeze rule that lets a start carry information through the iteration would be different designs, not tunings of this one, and the results here neither support nor exclude them. What the data suggest as a quantity to look for in a future design is the largest freeze index at which a given start still converges, and the share of geometries on which it does not. A start that moved that curve, one that converged on the geometries where a random start with an early freeze fails, would be a warm start in the sense of the proposal; the oracle start does it on three geometries of ten, and the network start does it on \capRNotN{} of the \nGeo{} held-out geometries while failing on \capNNotR{} where the random start with the same freeze converges. We also did not test whether a network trained at the published freeze index with a longer unroll would behave differently from the candidate; the single earlier trial at that setting reached about the random level, and it is the only evidence we have on the question. The oracle start suggests a saving of about one iteration for phase starts at the published freeze on these ten geometries, but it is one particular start (the converged phases of a single random-start run, one seed), so it is not a bound over all starts.
